\section{Version history}
\label{app:version-history}
\begingroup\small

Every numbered result of Version~1 keeps its number and place in later
versions, and its name except where noted below.

\subsection*{Version 3 (October 1, 2026)}

\begin{itemize}
  \item The paper was rewritten for readers outside the framework; it
    now says plainly that the conditional theorem is a reformulation of
    the Riemann hypothesis, not evidence for it.  Discussion of the formal
    development was moved to \Cref{app:formalization}; the numbered
    lists of non-claims and ``no-smuggling gates'' and the reference to
    a ``V-Differential'' classification, which is not defined in the
    cited literature, were removed; the account of a three-route
    derivation attempt for the recognition source was removed.
  \item \Cref{def:rh:nontrivial-zero-ledger} now allows arbitrary
    positive weights, with multiplicities as the natural choice;
    Version~2 required multiplicities, which the formal development does
    not use.  \Cref{thm:rh:translation-T} is stated for the actual
    nontrivial zeros and includes the equivalence with the Riemann
    hypothesis; the shell version is described in the text.
  \item \Cref{obl:rh:gamma-sdtc-selberg} is now stated in mathematical
    terms and displayed as a hypothesis (it was labelled
    ``Obligation'').  The new \Cref{prop:rh:gamma-calibration} states that for every shell it
    is equivalent to the critical-line statement for the represented
    zeros, hence on the bare shell to the Riemann hypothesis; Version~2
    stated the bare-shell case in prose.  \Cref{def:rh:sat-sel-shell}
    now states the coordinate system and zero ledger of a shell
    explicitly.
  \item The new \Cref{thm:rh:concrete-window-dc} applies the
    duality-confinement theorem, in the concrete operator form of the
    companion paper's Version~3, to finite reflection-closed windows of
    actual zeros that eventually contain every zero, and proves that the
    resulting domination records exist if and only if the Riemann
    hypothesis holds.
  \item The AOR classification (\Cref{sec:rh-aor-instance}) was
    condensed.  Items~8.2--8.8, which record classification labels and
    were displayed as theorems, are now displayed as ``Records'' and
    describe the recorded lists exactly; \Cref{thm:rh:aor-instance} now
    includes the non-emptiness of the non-claim list.  The section was
    renamed ``A classification of the hypotheses''.
  \item The finite-window material of the Version~2 discussion was
    reorganized into the new \Cref{sec:finite_windows}, with the new
    \Cref{prop:rh:window-identities,prop:rh:window-obstructions,prop:rh:window-equivalences}
    and complete proofs.
  \item Names shortened: \Cref{def:rh:sat-sel-shell} (formerly
    ``Saturated completed Selberg trace shell interface''),
    \Cref{thm:rh:dc-master-applied} (formerly ``Duality-confinement
    master theorem, applied to'' the shell),
    \Cref{thm:rh:full-classical} (formerly ``Full classical conditional
    RH''), Records~8.2--8.8 (formerly ``Mechanical AOR records'',
    ``Recognition-source discharge'', ``Bridge-field discharge'',
    ``Admissibility-audit status record'', ``Cross-paper master-theorem
    discharge'', ``Typed real-coordinate discharge'',
    ``Translation-interface discharge''), \Cref{thm:rh:aor-instance} (formerly ``RH local-register theorem''),
    and \Cref{rem:rh:aor-partial-status} (formerly ``Partial status of
    the AOR-instance wrapper'').
  \item Cross-references now name the kind of result they point to;
    bibliography entries were corrected.
\end{itemize}

\subsection*{Version 2 (September 28, 2026)}

Version~1 described the argument as a refinement-stable
audited-operational-realisability instance and identified the shell's
real-part coordinate with the classical real part by construction.
Version~2 replaced the first claim by the local-register statement,
which is what was proved; made the identification an explicit
hypothesis, constructed it for the bare shell, and proved the passage
from the critical strip to the full Riemann hypothesis; recorded that on
the bare shell the recognition source exists exactly under the Riemann
hypothesis; added the finite-window analysis and the comparison with
the companion Navier--Stokes and metatheory papers; and moved the formal
development from a library-free encoding to Mathlib, correcting the
description of what it proves.
\par\endgroup
